% GENERATED FILE. Edit canonical lessons, then run npm run generate:books.
% canonical-source-hash: fnv1a64:1b16ee5a9a6cb6ef
% canonical-lessons: JA-C435-manzoku, JA-C435-fuman, JA-C435-taikutsu, JA-C435-hohoemu, JA-C435-damaru

\chapter{Satisfied or Bored}
\label{ch:ja-satisfied-or-bored}
\hlchaptermodality{435}

\begin{chapteropening}
\textbf{By the end of this chapter:} \emph{I can say whether I am satisfied or bored.}

Complete the last lesson of chapter 435: I can say whether I am satisfied or bored.
\end{chapteropening}

\section[\texorpdfstring{\ja{まんぞく} (manzoku)}{まんぞく (manzoku)}]{\ja{まんぞく} (manzoku) — satisfied; contented}

\label{lesson:JA-C435-manzoku}



\begin{quote}
Before the new one: say the Japanese for miserable, then the Japanese for to encourage.
\end{quote}

\subsection*{You'll want to know: \ja{まんぞく}}
\textbf{\ja{まんぞく}} — \emph{manzoku} — \textquotedblleft{}satisfied; contented\textquotedblright{}.

\begin{quote}
\textbf{\ja{この} \ja{ばんごはんに} \ja{まんぞくしました。}} — \emph{kono bangohan ni manzoku shimashita} — \textquotedblleft{}I was satisfied with this dinner.\textquotedblright{}
\end{quote}

A na-adjective: \ja{な} before a noun.

\begin{cousinweb}
\textbf{\ja{まんぞく}} (\emph{manzoku}) is Sino-Japanese (kango): two Chinese characters read with their on'yomi, readings borrowed from Chinese. They are \emph{man}, \textquotedblleft{}full\textquotedblright{} (Middle Chinese \emph{manX}) and \emph{zoku}, \textquotedblleft{}be sufficient\textquotedblright{} (\emph{tsjowk}) (KANJIDIC2; Middle Chinese in Baxter's transcription).
\end{cousinweb}

\begin{grammarlens}[title={What you've built}]
One more word for describing what you lived and felt.
\end{grammarlens}

\subsection*{Your turn}
\begin{itemize}
\raggedright
  \item \emph{Say it:} \emph{manzoku}
  \item \emph{Say it:} \emph{manzoku}, once more
  \item \emph{Say it:} \emph{kono bangohan ni manzoku shimashita} — I was satisfied with this dinner.
  \item \emph{Recall:} say the Japanese for miserable, then the Japanese for to encourage, then say \emph{manzoku} again
\end{itemize}

\begin{tcolorbox}[breakable,colback=teal!4,colframe=teal!35!black,title={Before you move on}]
What does \ja{まんぞく} mean? (\textquotedblleft{}Satisfied; contented\textquotedblright{}.) Say it once more.
\end{tcolorbox}
\section[\texorpdfstring{\ja{ふまん} (fuman)}{ふまん (fuman)}]{\ja{ふまん} (fuman) — discontented; dissatisfied}

\label{lesson:JA-C435-fuman}



\begin{quote}
Before the new one: say the Japanese for to encourage, then the Japanese for satisfied.
\end{quote}

\subsection*{You'll want to know: \ja{ふまん}}
\textbf{\ja{ふまん}} — \emph{fuman} — \textquotedblleft{}discontented; dissatisfied\textquotedblright{}.

\begin{quote}
\textbf{\ja{あには} \ja{しごとに} \ja{ふまんが} \ja{あります。}} — \emph{ani wa shigoto ni fuman ga arimasu} — \textquotedblleft{}My elder brother is unhappy with his job.\textquotedblright{}
\end{quote}

A na-adjective: \ja{な} before a noun.

\begin{cousinweb}
\textbf{\ja{ふまん}} (\emph{fuman}) is Sino-Japanese (kango): two Chinese characters read with their on'yomi, readings borrowed from Chinese. They are \emph{fu}, \textquotedblleft{}negative\textquotedblright{} (Middle Chinese \emph{pjuw}) and \emph{man}, \textquotedblleft{}full\textquotedblright{} (\emph{manX}) (KANJIDIC2; Middle Chinese in Baxter's transcription).
\end{cousinweb}

\begin{grammarlens}[title={What you've built}]
One more word for describing what you lived and felt.
\end{grammarlens}

\subsection*{Your turn}
\begin{itemize}
\raggedright
  \item \emph{Say it:} \emph{fuman}
  \item \emph{Say it:} \emph{fuman}, once more
  \item \emph{Say it:} \emph{ani wa shigoto ni fuman ga arimasu} — My elder brother is unhappy with his job.
  \item \emph{Recall:} say the Japanese for to encourage, then the Japanese for satisfied, then say \emph{fuman} again
\end{itemize}

\begin{tcolorbox}[breakable,colback=teal!4,colframe=teal!35!black,title={Before you move on}]
What does \ja{ふまん} mean? (\textquotedblleft{}Discontented; dissatisfied\textquotedblright{}.) Say it once more.
\end{tcolorbox}
\section[\texorpdfstring{\ja{たいくつ} (taikutsu)}{たいくつ (taikutsu)}]{\ja{たいくつ} (taikutsu) — bored; tedious}

\label{lesson:JA-C435-taikutsu}



\begin{quote}
Before the new one: say the Japanese for satisfied, then the Japanese for discontented.
\end{quote}

\subsection*{You'll want to know: \ja{たいくつ}}
\textbf{\ja{たいくつ}} — \emph{taikutsu} — \textquotedblleft{}bored; tedious\textquotedblright{}.

\begin{quote}
\textbf{\ja{あめの} \ja{しゅうまつは} \ja{たいくつです。}} — \emph{ame no shūmatsu wa taikutsu desu} — \textquotedblleft{}A rainy weekend is boring.\textquotedblright{}
\end{quote}

A na-adjective: \ja{な} before a noun.

\begin{cousinweb}
\textbf{\ja{たいくつ}} (\emph{taikutsu}) is Sino-Japanese (kango): two Chinese characters read with their on'yomi, readings borrowed from Chinese. They are \emph{tai}, \textquotedblleft{}withdraw\textquotedblright{} (Middle Chinese \emph{thwojH}) and \emph{kutsu}, \textquotedblleft{}bend\textquotedblright{} (\emph{khjut}) (KANJIDIC2; Middle Chinese in Baxter's transcription).
\end{cousinweb}

\begin{grammarlens}[title={What you've built}]
One more word for describing what you lived and felt.
\end{grammarlens}

\subsection*{Your turn}
\begin{itemize}
\raggedright
  \item \emph{Say it:} \emph{taikutsu}
  \item \emph{Say it:} \emph{taikutsu}, once more
  \item \emph{Say it:} \emph{ame no shūmatsu wa taikutsu desu} — A rainy weekend is boring.
  \item \emph{Recall:} say the Japanese for satisfied, then the Japanese for discontented, then say \emph{taikutsu} again
\end{itemize}

\begin{tcolorbox}[breakable,colback=teal!4,colframe=teal!35!black,title={Before you move on}]
What does \ja{たいくつ} mean? (\textquotedblleft{}Bored; tedious\textquotedblright{}.) Say it once more.
\end{tcolorbox}
\section[\texorpdfstring{\ja{ほほえむ} (hohoemu)}{ほほえむ (hohoemu)}]{\ja{ほほえむ} (hohoemu) — to smile gently}

\label{lesson:JA-C435-hohoemu}



\begin{quote}
Before the new one: say the Japanese for discontented, then the Japanese for bored.
\end{quote}

\subsection*{You'll want to know: \ja{ほほえむ}}
\textbf{\ja{ほほえむ}} — \emph{hohoemu} — \textquotedblleft{}to smile gently\textquotedblright{}.

\begin{quote}
\textbf{\ja{かのじょは} \ja{すこし} \ja{ほほえみました。}} — \emph{kanojo wa sukoshi hohoemimashita} — \textquotedblleft{}She gave a little smile.\textquotedblright{}
\end{quote}

Polite \textbf{\ja{ほほえみます}}, \emph{hohoemimasu}; a godan verb. Past (polite): \textbf{\ja{そぼが} \ja{ほほえみました。}}, \emph{sobo ga hohoemimashita}.

\begin{cousinweb}
\textbf{\ja{ほほえむ}} (\emph{hohoemu}) is a native Japanese word (wago), inherited rather than borrowed from Chinese; its kanji spelling was attached to it later (JMdict).
\end{cousinweb}

\begin{grammarlens}[title={What you've built}]
One more word for describing what you lived and felt.
\end{grammarlens}

\subsection*{Your turn}
\begin{itemize}
\raggedright
  \item \emph{Say it:} \emph{hohoemu}
  \item \emph{Say it:} \emph{hohoemu}, once more
  \item \emph{Say it:} \emph{kanojo wa sukoshi hohoemimashita} — She gave a little smile.
  \item \emph{Recall:} say the Japanese for discontented, then the Japanese for bored, then say \emph{hohoemu} again
\end{itemize}

\begin{tcolorbox}[breakable,colback=teal!4,colframe=teal!35!black,title={Before you move on}]
What does \ja{ほほえむ} mean? (\textquotedblleft{}To smile gently\textquotedblright{}.) Say it once more.
\end{tcolorbox}
\section[\texorpdfstring{\ja{だまる} (damaru)}{だまる (damaru)}]{\ja{だまる} (damaru) — to fall silent}

\label{lesson:JA-C435-damaru}



\begin{quote}
Before the new one: say the Japanese for bored, then the Japanese for to smile gently.
\end{quote}

\subsection*{You'll want to know: \ja{だまる}}
\textbf{\ja{だまる}} — \emph{damaru} — \textquotedblleft{}to fall silent\textquotedblright{}.

\begin{quote}
\textbf{\ja{せんせいが} \ja{きて、みんな} \ja{だまりました。}} — \emph{sensei ga kite, minna damarimashita} — \textquotedblleft{}The teacher came, and everyone fell silent.\textquotedblright{}
\end{quote}

Polite \textbf{\ja{だまります}}, \emph{damarimasu}; a godan verb. Past (polite): \textbf{\ja{こどもが} \ja{だまりました。}}, \emph{kodomo ga damarimashita}.

\begin{cousinweb}
\textbf{\ja{だまる}} (\emph{damaru}) is a native Japanese word (wago), inherited rather than borrowed from Chinese; its kanji spelling was attached to it later (JMdict).
\end{cousinweb}

\begin{grammarlens}[title={What you've built}]
One more word for describing what you lived and felt.
\end{grammarlens}

\subsection*{Your turn}
\begin{itemize}
\raggedright
  \item \emph{Say it:} \emph{damaru}
  \item \emph{Say it:} \emph{damaru}, once more
  \item \emph{Say it:} \emph{sensei ga kite, minna damarimashita} — The teacher came, and everyone fell silent.
  \item \emph{Recall:} say the Japanese for bored, then the Japanese for to smile gently, then say \emph{damaru} again
\end{itemize}

\begin{tcolorbox}[breakable,colback=teal!4,colframe=teal!35!black,title={Before you move on}]
What does \ja{だまる} mean? (\textquotedblleft{}To fall silent\textquotedblright{}.) Say it once more.
\end{tcolorbox}
